\section{Récupération d'un modèle de RAM}
\label{sec:ram}

Le premier bloc à réutiliser est une RAM synchrone fournie avec son testbench.
Nous avons conservé son fonctionnement, puis adapté les types, les dimensions et la
sélection pour pouvoir relier plusieurs instances au même bus.

\subsection{Analyse du modèle et du testbench}

Le modèle initial, \texttt{ramchip.vhd}, décrit l'entité \texttt{sync\_ram}.
Son adresse de 10 bits donne accès à $2^{10}=1024$ mots de 8 bits. Les entrées sont
l'horloge \texttt{clock}, l'autorisation d'écriture \texttt{we}, l'adresse et la donnée à
écrire ; la sortie \texttt{dataout} fournit la donnée lue. Le tableau mémoire n'est pas
initialisé.

Le code~\ref{lst:ram-originale} distingue deux mécanismes. Au front montant, si
\texttt{we = '1'}, la donnée est écrite à l'adresse présente sur l'entrée. Le même
processus mémorise l'adresse dans \texttt{read\_address}, quelle que soit la valeur de
\texttt{we}. L'affectation concurrente de sortie utilise cette adresse mémorisée : un
changement de l'adresse d'entrée ne suffit donc pas à actualiser la lecture, il faut
attendre le front montant suivant.

\noindent\begin{minipage}{\linewidth}
\begin{lstlisting}[style=vhdl, caption={Lecture et écriture du modèle initial}, label={lst:ram-originale}]
if rising_edge(clock) then
  if we = '1' then
    ram(to_integer(unsigned(address))) <= datain;
  end if;
  read_address <= address;
end if;
...
dataout <= ram(to_integer(unsigned(read_address)));
\end{lstlisting}
\end{minipage}\par

Le fichier \texttt{ram\_tb.vhd} contient deux architectures : \texttt{Bench} présente
un scénario court et \texttt{Bench2} effectue une boucle d'écriture suivie d'une boucle de
lecture. Nous les avons simulées séparément. L'horloge a une période de 10~ns et s'arrête
lorsque \texttt{StopClock} devient vrai.

\paragraph{Scénario court.}
Le testbench modifie ses entrées juste après un front montant, par des affectations de
signaux. Le processus de la RAM voit encore les anciennes valeurs à ce front : les
nouvelles valeurs seront prises en compte au front suivant. Le
tableau~\ref{tab:ram-chrono} donne les principaux événements observés sur la
figure~\ref{fig:ram-originale}.

\begin{table}[H]
\centering
\begin{tabularx}{\linewidth}{rX}
\toprule
Temps & Résultat observé \\
\midrule
25~ns & Écriture de 4 au mot 1 ; la sortie prend la valeur 4 après mise à jour des signaux. \\
35~ns & Écriture de 7 au mot 2 ; la sortie prend la valeur 7. \\
45~ns & Le testbench demande le mot 1 ; l'adresse mémorisée reste 2 et la sortie reste 7. \\
55~ns & L'adresse 1 est mémorisée et la sortie prend la valeur 4. \\
65~ns & L'adresse 2 est mémorisée et la sortie prend la valeur 7. \\
\bottomrule
\end{tabularx}
\caption{Chronologie du test court de la RAM initiale}
\label{tab:ram-chrono}
\end{table}

\simcapture{p2-original.png}{615 564 10 130}{615 154 10 630}
{RAM initiale : écritures de 4 et 7, puis lectures aux adresses 1 et 2}
{fig:ram-originale}

La sortie initiale indéfinie est attendue : ni la mémoire ni le registre d'adresse ne
sont initialisés. ModelSim peut aussi signaler à 0~ns une métavaleur lors de la conversion
de l'adresse en entier. Enfin, quand un front écrit un mot et mémorise en même temps son
adresse, la sortie présente directement la nouvelle donnée : l'affectation concurrente de
sortie est réévaluée après la mise à jour du mot et de l'adresse.

\paragraph{Test en boucle.}
\texttt{Bench2} utilise le signal booléen \texttt{ok}, qui reste faux dès qu'une comparaison
échoue. Notre exécution se termine à 20490~ns avec \texttt{ok = TRUE}. La comparaison
soustrait 1 à l'adresse courante car, au moment où elle est exécutée au front montant,
\texttt{dataout} correspond encore à l'adresse précédente.

Le test a toutefois deux limites : les données se répètent tous les 256 mots, ce qui peut
masquer un défaut sur les bits d'adresse de poids fort, et la dernière donnée lue au mot
1023 n'est pas comparée. Nous avons gardé ce test pour la comparaison avant/après, puis
ajouté un contrôle de toutes les adresses à la section~\ref{sec:ram-generique}.

\subsection{Passage au type unsigned et nouvelle simulation}

Nous avons remplacé \texttt{std\_logic\_vector} par \texttt{unsigned} pour les ports
d'adresse et de données, les mots mémoire, le registre d'adresse et les signaux du
testbench. Ce type de \texttt{numeric\_std} représente un entier non signé et permet
d'écrire directement \texttt{address + 1}.

Les conversions entourant l'adresse deviennent inutiles :
\texttt{to\_integer(unsigned(address))} devient \texttt{to\_integer(address)}.
Le rôle des signaux, leurs largeurs et les conditions de lecture et d'écriture sont
conservés. L'horloge et \texttt{we}, qui ne portent pas des nombres, restent de type
\texttt{std\_logic}.

Nous avons relancé les deux architectures après recompilation. Le scénario court
reproduit les valeurs et les instants du tableau~\ref{tab:ram-chrono}
(figure~\ref{fig:ram-unsigned}). Le test \texttt{Bench2} se termine à nouveau à 20490~ns,
avec \texttt{ok = TRUE}. Ces deux observations confirment la conservation du comportement
pour les scénarios testés.

\simcapture{p2-unsigned.png}{660 564 10 130}{660 154 10 630}
{RAM avec ports unsigned : même chronologie que le modèle initial}
{fig:ram-unsigned}

\subsection{Paramétrage de la capacité et de la largeur des données}
\label{sec:ram-generique}

Deux paramètres \texttt{generic} déterminent les dimensions du bloc
(code~\ref{lst:ram-generique}). \texttt{Nbits\_addr} donne la largeur d'adresse, donc la
capacité vaut $2^{\mathtt{Nbits\_addr}}$ mots ; \texttt{Nbits\_data} donne le nombre de bits
par mot. Nous les avons déclarés de type \texttt{positive} pour exclure les largeurs nulles.
Les valeurs par défaut, 10 et 8, retrouvent le modèle initial.

\noindent\begin{minipage}{\linewidth}
\begin{lstlisting}[style=vhdl, caption={Paramètres et ports de la RAM générique}, label={lst:ram-generique}]
generic (
  Nbits_addr : positive := 10;
  Nbits_data : positive := 8
);
port (
  clock   : in  std_logic;
  we      : in  std_logic;
  address : in  unsigned(Nbits_addr-1 downto 0);
  datain  : in  unsigned(Nbits_data-1 downto 0);
  dataout : out unsigned(Nbits_data-1 downto 0)
);
\end{lstlisting}
\end{minipage}\par

L'architecture utilise déjà les attributs des ports : le tableau contient
\texttt{2**address'length} éléments, chacun de type \texttt{unsigned(datain'range)}, et
\texttt{read\_address} reprend \texttt{address'range}. Elle s'adapte ainsi aux dimensions
choisies sans changer le processus de lecture et d'écriture. Un \texttt{generic map}
permet de choisir une configuration propre à chaque instance.

\paragraph{Vérification de deux configurations.}
Le \texttt{Bench2} hérité compare encore les données à
\texttt{address(7 downto 0) - 1} : il n'est pas adapté à toutes les nouvelles dimensions.
Nous avons donc ajouté l'architecture \texttt{GenericCheck} dans
\path{ram_generic_check.vhd}, en conservant les deux tests précédents.

Ce test écrit puis relit toutes les adresses, y compris la dernière. Les entrées changent
au front descendant, la RAM les échantillonne au front montant et la lecture est contrôlée
au front descendant suivant. Le motif de données replie les bits d'adresse par XOR sur
la largeur du mot ; une seconde passe utilise son complément. Chaque bit de données est
ainsi exercé à 0 et à 1, et tous les bits d'adresse participent au motif. Les deux
configurations du tableau~\ref{tab:ram-generiques} passent ce contrôle.

\begin{table}[H]
\centering
\begin{tabular}{rrrrl}
\toprule
Bits d'adresse & Bits de données & Mots & Fin de simulation & Résultat \\
\midrule
4 & 16 & 16 & 660~ns & \texttt{ok = TRUE} \\
10 & 32 & 1024 & 40980~ns & \texttt{ok = TRUE} \\
\bottomrule
\end{tabular}
\caption{Configurations vérifiées avec deux passes complémentaires}
\label{tab:ram-generiques}
\end{table}

Le modèle de RAM et le testbench sont identiques pour les deux essais ; seuls les
paramètres changent au chargement dans ModelSim. Le contrôle se termine respectivement
à 650 et 40970~ns, puis l'horloge s'arrête une période plus tard. La
figure~\ref{fig:ram-g32} montre les deux passes pour la configuration 1024 mots de 32 bits.

\simcapture{p2-generic-1024x32.png}{660 574 10 130}{660 154 10 630}
{Configuration 1024 mots de 32 bits : même testbench et résultat TRUE}
{fig:ram-g32}

\subsection{Ajout d'une sélection asynchrone CS}

Pour relier plusieurs sorties mémoire, une seule RAM doit piloter le bus à un instant
donné. Nous avons ajouté l'entrée \texttt{cs}, active à 1. Une
affectation concurrente impose la haute impédance lorsque la RAM n'est pas sélectionnée
(code~\ref{lst:ram-cs}). La réaction à \texttt{cs} est asynchrone : elle n'attend pas
de front d'horloge. L'adresse de lecture reste, elle, mémorisée au front montant.

\noindent\begin{minipage}{\linewidth}
\begin{lstlisting}[style=vhdl, caption={Sélection des écritures et commande de la sortie}, label={lst:ram-cs}]
if rising_edge(clock) then
  if we = '1' and cs = '1' then
    ram(to_integer(address)) <= datain;
  end if;
  read_address <= address;
end if;
...
dataout <= ram(to_integer(read_address))
           when cs = '1' else (others => 'Z');
\end{lstlisting}
\end{minipage}\par

La sélection doit aussi conditionner l'écriture : avec un \texttt{we} commun, les RAM
non sélectionnées seraient sinon modifiées. Les essais de la RAM avec sélection et de
\texttt{quadram} ont été relancés après la correction décrite en
section~\ref{sec:validsysteme}.

Le scénario court conserve d'abord les écritures de 4 et 7, puis désactive \texttt{cs}
à 67~ns, entre deux fronts d'horloge. La sortie passe immédiatement en haute impédance.
Au front montant de 85~ns, \texttt{we} vaut 1 et la donnée proposée est 9 à l'adresse 1,
mais \texttt{cs} vaut 0 : le mot reste égal à 4. À 87~ns, la sélection est rétablie,
toujours entre deux fronts, et la sortie reprend immédiatement la valeur 4
(figure~\ref{fig:ram-cs}). La trace interne \texttt{ram(1)} confirme la conservation du mot.

\simcapture{p2-cs.png}{660 564 10 130}{660 154 10 630}
{CS asynchrone : sortie à haute impédance et écriture bloquée entre 67 et 87 ns}
{fig:ram-cs}
